\documentclass[../main.tex]{subfiles}

\begin{document}

\chapter{ Week 2 }

代靖涵 2512022201319

\begin{exercise}{}{}
  证明: 度量空间 $X$ 中的序列 $(x_n)$ 收敛于 $p$ 当且仅当对于任何包含 $p$ 的开集 $U$, 总存在 $N \in \mathbb{Z}_+$, 使得当 $n \ge N$ 时有 $x_n \in U$.
\end{exercise}
\begin{proof}
    \begin{enumerate}
        \item \(  \implies   \): 若 \(  x_{n}\to p  \). 任取包含 \(  p  \)的开集 \(  U  \), 存在 \(   \delta > 0  \), 使得 \(  p \in N_{d}\left( p, \delta  \right)\subseteq U   \). 由收敛的定义, 存在 \(  N\in \mathbb{Z} _{+ }  \), 使得当 \(  n\ge N  \)时, 都有 \(  d\left( x_{n},p \right)<  \delta    \), 此时 \(  x_{n}\in N_{d}\left( p, \delta  \right)\subseteq U   \).    
        \item \(  \impliedby   \): 若对于任何包含 \(  p  \)的开集, 总存在 \(  N\in \mathbb{Z} _{+ }  \), 使得当 \(  n\ge N  \)时有 \(  x_{n}\in U  \). 特别地, 任取定 \(   \varepsilon > 0  \),  取 \(  U= N_{d}\left( p,  \varepsilon  \right)    \). 则 \(  x_{n}\in U  \)即 \(  d\left( x_{n},p \right)<  \varepsilon    \). 由于 \(   \varepsilon > 0  \)是任取的,   这就根据定义证明了 \(  \left( x_{n} \right)   \)收敛于 \(  p  \).          
    \end{enumerate}
    
\end{proof}


\begin{exercise}{}{}
    
我们用下面的反例说明集合 $W = \{f : [0,1] \to \mathbb{R} \mid f \text{ 是连续函数}\}$ 在度量 $d(f,g) = \int_0^1 |f(x) - g(x)|dx$ 下  是不完备的. 令 $f_n : [0,1] \to \mathbb{R}$,
$$f_n(x) = \begin{cases}
1, & 0 \le x \le \frac{1}{2}, \\
1 - 2^n(x - \frac{1}{2}), & \frac{1}{2} < x < \frac{1}{2} + (\frac{1}{2})^n, \\
0, & \frac{1}{2} + (\frac{1}{2})^n \le x \le 1.
\end{cases}$$
证明 $(f_n)$ 是 $(W,d)$ 中的 Cauchy 序列, 但不在 $W$ 中收敛.
\end{exercise}

\begin{proof}
    任取 \(   \varepsilon > 0  \), 存在 \(  N  \), 使得\(  \frac{1 }{2^{N+ 1} }<  \varepsilon    \).  则对于任意的 \(  m\ge n  \ge N\),都 有  \[
  \begin{aligned}
  \int_{0}^{1}\left| f_{m}-f_{n} \right|\,\mathrm{d} x&=   \int_{\frac{1}{2}}^{\frac{1}{2}+ \frac{1 }{2^{m} } }  \left| f_{m}-f_{n} \right|\,\mathrm{d} x + \int_{\frac{1}{2}+ \frac{1 }{2^{m} } }^{\frac{1}{2}+ \frac{1 }{2^{n} } }\left| f_{n} \right|\,\mathrm{d} x\\ 
   &= \int_{\frac{1}{2}}^{\frac{1}{2}+ \frac{1 }{2^{m} } } \left( x-\frac{1 }{2 }  \right)\left( 2^{m}-2^{n} \right)  \,\mathrm{d} x+  \int_{\frac{1}{2}+ \frac{1 }{2^{m} } }^{\frac{1}{2}+ \frac{1 }{2^{n} } }\left( 1-2^{n}\left( x-\frac{1 }{2 }  \right)  \right)\,\mathrm{d} x\\ 
    &= \left( 2^{m}-2^{n} \right)\frac{1}{2}\left( \frac{1 }{2^{2m} }  \right)+ \left( \frac{1 }{2^{n} }-\frac{1 }{2^{m} }   \right)- 2^{n} \frac{1 }{2 }\left( \frac{1 }{2^{2n} }-\frac{1 }{2^{2m} }   \right)      \\ 
     &= \frac{1}{2}\frac{1 }{2^{m} } -\frac{1}{2}\frac{1 }{2^{2m-n} } + \frac{1 }{2^{n} }-\frac{1 }{2^{m} } -\frac{1}{2}\frac{1 }{2^{n} }+ \frac{1}{2}\frac{1 }{2^{2m-n} }  \\ 
      &= \frac{1 }{2^{n+ 1} }-\frac{1 }{2^{m+ 1} }\\ 
       &\le \frac{1 }{2^{n+ 1} }\le \frac{1 }{2^{N+ 1} }<  \varepsilon     
  \end{aligned}
    \] 这表明 \(  \left( f_{n} \right)   \)是  \(  \left( W,d \right)   \)中的Cauchy列.  但是若 \(  f_{n}  \)收敛于 连续函数\(  f  \), 则 \[
    \int_{0}^{1}\left| f_{n}\left( x \right)-f\left( x \right)   \right|\,\mathrm{d} x\to 0 
    \]    \[
    \begin{aligned}
    &\int_{0}^{1}\left| f_{n}\left( x \right)-f\left( x \right)   \right|\,\mathrm{d} x  \\ 
     &=  \int_{0}^{\frac{1}{2}}\left| 1-f\left( x \right)  \right|\,\mathrm{d} x+  \int_{\frac{1}{2}}^{\frac{1}{2}+ \frac{1 }{2^{n} } }\left| 1-2^{n}\left( x-\frac{1 }{2 }  \right)-f\left( x \right)   \right|  +\int_{\frac{1 }{2 }+  \frac{1}{2^{n}}}^{1}\left| f\left( x \right)  \right|\,\mathrm{d} x 
    \end{aligned}
    \] 这三项当 \(  n\to \infty  \)时, 分别趋于零.  其中, 由于第一项为常数, 我们有 \[
    \int_{0}^{\frac{1}{2}}\left| 1-f\left( x \right)  \right|\,\mathrm{d} x= 0\implies  f\left( x \right)= 1,\forall x\in \left[ 0,\frac{1 }{2 }  \right]   
    \] 存在 \(  m   \), 使得当 \(  x \in \left[ \frac{1}{2},\frac{1 }{2 }+ \frac{1 }{2^{m} }   \right]   \)时, 都有 \(  \left| f\left( x \right)-1  \right|< \frac{1 }{2 }    \)   此时 对于任意的 \(  n\ge m+ 1  \), 都有 \[
    \int_{\frac{1}{2}+ \frac{1 }{2^{n} } }^{1}\left| f\left( x \right)  \right|\ge  \int_{\frac{1}{2}+ \frac{1 }{2^{n} } }^{\frac{1}{2}+ \frac{1 }{2^{m} }}\left( 1-\left| f\left( x \right)-1  \right|  \right) \ge \frac{1}{2}\left( \frac{1 }{2^{m} }-\frac{1 }{2^{n} }   \right) \ge \frac{1 }{2^{m+ 2} } 
    \]这与 \[
    \int_{\frac{1}{2}+ \frac{1 }{2^{n} } }^{1}\left| f\left( x \right)  \right|\to 0,\quad (n\to \infty) 
    \]矛盾. 因此不存在连续函数 \(  f  \), 使得 \(  \left( f_{n} \right)   \)收敛与 \(  f  \), 即 \(  \left( f_{n} \right)   \)不在 \(  W  \)中收敛.     
\end{proof}

\begin{exercise}{}{}
  证明以下与 \(  \mathbb{R}   \)  有关的运算是连续的
  \begin{enumerate}
    \item 加法: $\mathbb{R} \times \mathbb{R} \to \mathbb{R}, (a, b) \mapsto a + b$.
    \item 取相反数: $\mathbb{R} \to \mathbb{R}, a \mapsto -a$.
    \item 乘法: $\mathbb{R} \times \mathbb{R} \to \mathbb{R}, (a, b) \mapsto ab$.
    \item 取倒数: $\mathbb{R} - \{0\} \to \mathbb{R} - \{0\}, a \mapsto a^{-1}$.
    \item 取最大值: $\mathbb{R} \times \mathbb{R} \to \mathbb{R}, (a, b) \mapsto \max\{a, b\}$.
    \item 取最小值: $\mathbb{R} \times \mathbb{R} \to \mathbb{R}, (a, b) \mapsto \min\{a, b\}$.
\end{enumerate}
\end{exercise}

\begin{proof}
  \begin{enumerate}
    \item 对于任意的 \(  \left( a,b \right)\in \mathbb{R} ^{2}   \),任取 \(   \varepsilon > 0  \). 我们有\(   \forall \left( x,y \right) \in B_{d\times d}\left( \left( a,b \right),\frac{ \varepsilon }{2 }   \right)  \), 都有 \(  \left| x-a \right|< \frac{ \varepsilon  }{2 },  \left| y-b \right|< \frac{ \varepsilon  }{2 }       \) , 从而 \[
    \left| \left( x+ y \right)-\left( a+ b \right)   \right|\le \left| x-a \right|+ \left| y-b \right|<  \varepsilon     
    \]这表明 \(  x+ y \in B_{d}\left( a+ b, \varepsilon  \right)   \) . 从而加法运算是连续的.
    \item 对于任意的 \(  a \in \mathbb{R}  \),任取 \(   \varepsilon > 0  \), 则对于任意的 \(  x \in B_{d}\left( a,  \varepsilon  \right)   \), 都有\[
    \left| -x-\left( -a \right)  \right|= \left| x-a \right|<  \varepsilon   
    \], 因此 \[
    -x \in B_{d}\left( -a, \varepsilon  \right) 
    \]这映射 \(  a\mapsto -a  \)连续.
    \item 对于任意的 \(  \left( a,b \right)\in \mathbb{R} ^{2}   \),任取 \(   \varepsilon > 0  \). 取 \[
     \delta = \min \left\{ 1, \frac{ \varepsilon  }{\left| a \right|+ \left| b \right|+ 1   }  \right\}
    \]对于任意的 \(  \left( x,y \right)\in B_{d\times d}\left( \left( a,b \right), \delta   \right)     \), 都有 \(  \left| x-a \right|<  \delta    \), \(  \left| y-b \right|<  \delta    \), 进而   \[
    \begin{aligned}
    \left| xy-ab \right|=  \left| xy-ya+ ya-ab \right|&\le \left| y \right|\left| x-a \right|+ \left| a \right|\left| y-b \right|    \\ 
     &<  \left( \left| b \right|+  \delta   \right) \delta + \left| a \right| \delta \\ 
      & \le  \left( \left| a \right|+ \left| b \right|+ 1   \right)   \delta \\ 
       &\le  \varepsilon 
    \end{aligned}
    \]从而 \[
    xy \in B_{d}\left( ab,  \varepsilon  \right) 
    \]
    \item  对于任意的 \( a \in \mathbb{R} \setminus \left\{ 0 \right\}\), 任取 \(   \varepsilon > 0  \),  取 \[
     \delta = \min \left\{ \frac{1 }{2 }\left| a \right|, \frac{\left| a \right|^{2}   \varepsilon }{2 }    \right\}
    \]则对于任意的 \(  x\in B_{d}\left( a ,  \delta   \right)   \),  \[
    \left| x^{-1} -a^{-1}  \right|= \frac{\left| x-a \right|  }{\left| x \right|\left| a \right|   }  < \frac{ \delta  }{\left( \left| a \right|-\left| x-a \right|   \right) \left| a \right|  }\le \frac{ 2\delta  }{\left| a \right|^{2}   }  \le  \varepsilon 
    \]从而 \(  x^{-1} \in B_{d}\left( a^{-1} ,  \varepsilon  \right)   \), 这表明映射连续.
    \item 注意到 \[
    \max \left\{ a,b \right\}= \frac{a+ b+ \left| a-b \right|  }{2 } 
    \]由于加法, 加法逆 和数乘运算都是连续的, 且连续映射的复合映射连续. 因此只需要说明映射 \(  a\to \left| a \right|   \)是连续的. 任取 \(   \varepsilon > 0  \), 连续性由 \[
    \left| \left| a \right|-\left| x \right|   \right|\le \left| a-x \right|  
    \]立即得到 
    \item 注意到 \[
    \min \left\{ a,b \right\}= -\max \left\{ -a,-b \right\}
    \]由2. 5.可知6.中写作连续映射的复合映射, 从而也是连续映射 
  \end{enumerate}
  
   
\end{proof}

\begin{exercise}{}{}
利用推论 6.1.4 证明下面的集合是 $\mathbb{R}^2$ 的闭集:
\begin{enumerate}
\item $\{(x, y) | xy = 1\}$.
\item $\{(x, y) | x^2 + y^2 = 1\}$.
\item $\{(x, y) | x^2 + y^2 \le 1\}$.
\end{enumerate}
\end{exercise}

\begin{proof}
  \begin{enumerate}
    \item \(  \left( x,y \right)\mapsto xy   \)是 \(  \mathbb{R} ^{2}  \)  到 \(  \mathbb{R}   \)的连续映射. 集合\(  \left\{ 1 \right\}  \)不存在 \(  \mathbb{R}   \)上的 极限点, 是 \(  \mathbb{R}   \)的闭集.   故集合 \[
    \left\{ \left( x,y \right):xy= 1  \right\}
    \]是闭集\(  \left\{ 1 \right\}  \)在一个连续映射下的原像, 从而也是一个闭集.
    \item 考虑映射 \(  \mathbb{R} ^{2}\to \mathbb{R} ^{2}\to \mathbb{R}   \),  \(  \left( x,y \right)\mapsto   \left( x^{2},y^{2} \right)\mapsto x^{2}+ y^{2}  \)  . 由于连续映射的复合映射和乘积映射都是连续映射, 从而该映射也是连续的. 因此集合\[
    \left\{ \left( x,y \right):x^{2}+ y^{2}= 1  \right\}
    \]是 闭集\(  \left\{ 1 \right\}  \)在连续映射下的原像, 从而也是一个闭集. 
    \item  集合 \(  (-\infty,1]  \)作为开区间\(  (1,\infty)  \)的补集, 是一个闭集. 2.说明了映射 \(  \left( x,y \right)\to \left( x^{2}+ y^{2} \right)    \)是连续映射.  故集合 \[
    \left\{ \left( x,y \right):x^{2}+ y^{2}\le 1  \right\}
    \] 做为闭集 \(  (-\infty,1]  \) 在一个连续映射下的原像, 是一个闭集.
  \end{enumerate}
  
\end{proof}

\begin{exercise}{}{9-13-1}
证明不存在从闭区间 $[a, b]$ 到开区间 $(a, b)$ 上的连续满射.
\end{exercise}

\begin{proof}
  若 \(  f: [a,b]  \to (a,b)\)是一个连续的满射. 考虑 集合列 \(  \left\{ \left( a, b-\frac{1 }{n }  \right)  \right\}_{ n \in \mathbb{N} }  \). 构成 \(  \left( a,b \right)   \)的一个开覆盖, \[
  \left( a,b \right)= \bigcup _{n = 1}^{\infty}\left( a,b-\frac{1 }{n }  \right)  
  \]  这里规定 \(  \left( c,d \right)\coloneqq \varnothing   \)若 \(  d\le  c  \).   我们有 \[
  \left[ a,b \right]= f^{-1} \left( \left( a,b \right)  \right)= \bigcup _{n = 1}^{\infty}f^{-1} \left( \left( a, b-\frac{1 }{n }  \right)  \right)   
  \] 由于 \(  f  \)是连续映射, 每个 \(  f^{-1} \left( \left( a,b-\frac{1 }{n }  \right)  \right)   \)都是一个开集.  由此 \(  \left\{ f^{-1} \left( \left( a,b-\frac{1 }{n }  \right)  \right)  \right\}  \)构成 \(  [a,b]  \)的一个开覆盖. 从而存在有限子覆盖 \(  \left\{ f^{-1} \left( \left( a,b-\frac{1 }{n_1 }  \right)  \right),\cdots , f ^{-1} \left( \left( a,b-\frac{1 }{n_{k} }  \right)  \right)   \right\}  \)   . 不妨设 \(  n_{k}\ge n_{k-1}\ge \cdots \ge n_{1}  \). 则 \[
  [a,b]\subseteq f^{-1} \left( \left( a,b-\frac{1 }{n_{k} }  \right)  \right) 
  \] 由于 \(  f  \)是满射, 故 \[
  \left( a,b \right)= f\left( [a,b] \right)=  f \left( f^{-1} \left( a,b- \frac{1 }{n_{k} }  \right)  \right)= \left( a, b-\frac{1 }{n_{k} }  \right)    ,
  \]  矛盾. 因此不存在这样的连续满射.
\end{proof}
\begin{exercise}{}{}
证明 $\mathbb{R}$ 与开区间 $(0, 1)$ 是同胚的. $\mathbb{R}$ 与闭区间 $[0, 1]$ 是否同胚?
\end{exercise}

\begin{proof}
  考虑复合映射 \(  (0,1)\to (-\frac{ \pi }{2 },\frac{ \pi }{2 }  ) \to \mathbb{R}  \)   \[
  x\mapsto \frac{ \pi }{2 }\left( x-\frac{1 }{2 }  \right) \to \tan \left( \frac{ \pi }{2 }\left( x-\frac{1 }{2 }  \right)   \right)  
  \]是同胚映射 \(  (0,1)\to \left( -\frac{ \pi }{2 },\frac{ \pi }{2 }   \right)   \), \(  x\mapsto \frac{ \pi }{2 }\left( x-\frac{1 }{2 }  \right)    \)和另一个同胚映射 \(  \left( -\frac{ \pi }{2 },\frac{ \pi }{2 }   \right)\to \mathbb{R}    \),  \(  y\mapsto  \tan y  \)的复合映射.  从而是一个同胚映射. 因此 \(  \mathbb{R}   \)与 \(  (0,1)  \)同胚.    Exercise \ref{ex:9-13-1} 说明了不存在 \(   \left[ 0,1 \right]   \)到 \(  (0,1) \)的连续满射, 特别地, 不存在 \(  \left[ 0,1 \right]  \)到 \(   (0,1) \)的同胚映射.    因此 \(  (0,1)  \)与 \(  \left[ 0,1 \right]   \)不同胚. 而同胚是拓扑空间全体上的一个等价关系. 因此\(  \mathbb{R}   \)与 \(  \left[ 0,1 \right]   \)不同胚.    
\end{proof}


\begin{exercise}{}{}
 设 $f:(X,d) \to (Y,\rho)$ 是映射. 如果对于 $X$ 中所有的开 (闭) 集 $U$, $f(U)$ 都是 $Y$ 的开 (闭) 集, 则称 $f$ 是 \dfntxt{开 (闭) 映射}. 证明如果 $f$ 同时是连续的双射和开 (闭) 映射, 则 $f$ 是同胚.
\end{exercise}

\begin{proof}
  只证明开映射的命题, 闭映射的命题是完全类似地. 若 \(  f  \)连续而开的双射, 证明 \(  f  \)同胚只需要说明它的逆映射\(  F\left( x \right)\coloneqq f^{-1} \left( x \right)    \)     也是连续映射. 任取\(  X  \)上的开集 \(  U  \), 由于 \(  f  \)是开映射,  \(  V\coloneqq f\left( U \right)   \)也是一个开集. 由于 \(  F  \)是\(  f  \)的逆映射 ,  \[
  F^{-1} \left( U \right)= f\left( U \right) = V
  \]是一个开集. 由于 \(  U  \)是任取的, 因此 \(  F  \)也是一个连续映射. 
\end{proof}


\end{document}